\documentclass[12pt]{article}

\usepackage{amsmath}
\usepackage{german}
\usepackage{graphicx}
\DeclareMathOperator\sgn{sgn}

%\graphicspath{{/home/ritz/Teaching/fcn/fig/}}

%----------------- layout ----------------------------------------

\oddsidemargin=0in
\textwidth=6in
\topmargin=-0.7in
\textheight=9.5in
\baselineskip=16pt
\font\mittel=cmr17
\font\roman=cmr12
\font\klein=cmr9
\font\mgross=cmr10
\newfont{\CAP}{cmcsc10 scaled 900}
\pagestyle{empty}
\parindent=0mm

%----------------- macros ----------------------------------------
\newcommand{\conc}[1]{[{\rm C}]_{\text{#1}}}
\newcommand{\units}[1]{\text{#1}}
\newcommand{\K}{\text{K}}
\newcommand{\Na}{\text{Na}}
\newcommand{\Cl}{\text{Cl}}
\newcommand{\Leak}{\text{leak}}
\newcommand{\water}{\text{H}_2\text{O}}
\newcommand{\tderiv}[1]{\frac{\mathrm{d}#1}{\mathrm{d}t}}
%-----------------------------------------------------------------


\begin{document}

\parbox{2cm}{
\includegraphics[width=1.8cm]{bccnlogo.pdf}
}
\parbox{11cm}{
\begin{center}
\large HUMBOLDT-UNIVERSIT"AT \hskip 0.1 cm ZU \hskip 0.1 cm BERLIN
\vskip 0.1 true cm
\mgross BERNSTEIN CENTER FOR\\ COMPUTATIONAL NEUROSCIENCE
\end{center}
}
\parbox{2cm}
{
\hfill
\includegraphics[width=1.8cm]{hublogo.pdf}
}

\vskip 0.6 true cm

\leftline{\CAP Humboldt-Universit"at zu Berlin 
\hfill Phone: 030/2093-9110}
\leftline{\CAP Philippstr. 13 House 6
              \hfill Fax: 030/2093-6771}
\leftline{\CAP \hfill webpage: http://www.bccn-berlin.de/}


\vskip 0.6 true cm
\centerline{\bf Models of Neural Systems, WS 2010/11}
\centerline{\bf Project 7: FitzHugh-Nagumo model}
\centerline{Project presentation and report submission: February, 15th, 2011}

\vskip 0.8 true cm

\noindent {\bf Background}

Many features of dynamical systems can be understood by geometrical
analysis of the underlaying equations. This method is called phase
space analysis. Here we will considere a simplified two-dimensional
system of spike generation proposed by Richard FitzHugh:

\begin{eqnarray} 
\dot{V} & = & V-V^3/3 - W + I \\ 
\dot{W} & = & 0.08(V+0.7 - 0.8W)
\end{eqnarray} 

\vskip 0.2 true cm

{\bf Problems}


\begin{enumerate}
    \item \textbf{Phase space}
        \begin{itemize}
        \item Plot the vector flow field of the FitzHugh-Nagumo system on
            a grid $V=$-2.5 -- 2.5, $W=$-1.25 -- 1.25. Take $I=0$.
            \textit{Hint}: You can generate a 2D grid of points with
            \texttt{pylab.meshgrid} function; a field of arrows can be
            plotted in Pylab with \texttt{pylab.quiver} (check
            documentation for arguments).  
        \item Draw the $W$, $V$ nullclines on the same plot. What does
            the vector field look like near the nullclines and their
            intersection? 
        \item \textit{Stable points}. Starting at different initial
            conditions show that all trajectories go to the same
            point. Which point is it? Which initial conditions lead to
            the generation of an action potential? Can you find the
            potential ($V$) threshold above which spikes are reliably
            produced? \textit{Hint}: You can read out the coordinates
            of your mouse cursor in the lower-right corner of the plot.
        \item \textit{Limit cycle}. Increase external current to
                $I=1$ and plot several trajectories with chosen initial conditions. What has happened?
        \item \textit{Excitation block}. The FitzHugh-Nagumo model
            explains the excitation block phenomenon, i.e., the
            cessation of repetitive spiking as the amplitude of the
            stimulus current increases. Plot sample trajectories for
            $I=2$ and compare them with trajectories from the previous
            exercise. How can you explain the excitation block phenomenon?
        \item \textit{Post-inhibitory spikes}. In Hodgkin-Huxley and
            Fitzhugh-Nagumo release from inhibition can also lead to a
            spike, called post-inhibitory spike. In order to simulate
            this phenomenon find (numerically or geometrically) the
            stable fixed point for inhibitory input ($I$=-0.5). Next,
            use the values as a (forced) initial condition in a model
            without external current ($I$=0) and show that it leads to
            generation of an action potential.

       \end{itemize}
   \item \textbf{Bifurcation diagram}

       In this exercise you will plot bifurcation diagram of
       FitzHugh-Nagumo model, which summarizes its solutions as the
       bifurcation parameter (external current, $I$) is varied.
       \begin{itemize}
    \item \textit{Fixed points} Find fixed points of Fitzugh-Nagumo
        model as a function of $I$ by numerically
        solving the equation $V^3/3-V+(V+0.7)/0.8-I=0$. Take
        $I=$-0.5--2.5. Plot the value of the membrane potential ($V$)
        at the fixed point as a function of $I$.
        \textit{Hint:} Use \texttt{scipy.optimize.newton(func, x0)}
        which finds zeros of a function defined by \texttt{func}. Take
        \texttt{x0 = 1}.
    \item \textit{Stability}. Determine stability of the fixed points
        numerically. Simulate the model with an initial condition
        very close to the fixed point (but not in the fixed point) and
        check if the trajectory stays in its proximity (stable point)
        or diverges (unstable point). Repeat the procedure for each
        fixed point. Plot the stable points with a solid line and
        unstable points with a dashed line.
    \item \textit{Stable limit cycles}. Determine the position of
        stable limit cycles. Start the integration from the initial
        condition $(V(0), W(0))=(-10,10)$. Plot the minimum and
        maximum of the trajectory as a function of $I$.  
    
    \item \textit{Unstable limit cycles}. Determine the
        position of unstable limit cycles by repeating the last step
        but with integration time reversed ($t\rightarrow-t$). Use the
        initial condition $(V(0), W(0))=(1.,1.)$.
    
    \item Show all of the results on a bifurcation diagram.  

    \item Check what happens at the bifurcation point by narrowing the
        bifurcation parameter range (e.g. $I=0.32-0.34$). What
        type of bifurcation is it?
    \end{itemize}
\end{enumerate}

\begin{thebibliography}{7}
    \bibitem{dayan2001} P Dayan, LF Abbott, ``Theoretical neuroscience'',
        MIT Press 2001
    \bibitem{white2000} JA White, JT Rubinstein, AR Kay, ``Channel
        noise in neurons'', Trends Neurosci. 23: 131-137 (2000)
    \bibitem{chow1996} CC Chow, JA White, ``Spontaneous action
      potentials due to channel fluctuations'', Biophysical Journal
      71: 3013-3021 (1996)
\end{thebibliography}

\vfill
\centerline{\CAP Contact}
\CAP

\begin{tabular}{lll}
\underline{Matthias Guggenmos} & Phone: 450-517131 & Email: matthias.guggenmos(at)bccn-berlin.de \\
Urs Bergmann & Phone: 2093-8924 & Email: urs.bergmann(at)biologie.hu-berlin.de \\
Bartosz Telenczuk & Phone: 2093-8838 & Email: b.telenczuk(at)biologie.hu-berlin.de \\
Richard Kempter & Phone: 2093-8925 & Email: r.kempter(at)biologie.hu-berlin.de \\
\end{tabular}
\end{document}

 
